\documentclass{standalone}
\usepackage{circuitikz}
\usetikzlibrary{calc}
\definecolor{circuitnotemark}{HTML}{FE00FE}
\makeatletter
\pgfdeclareshape{smacoax}{
  \anchor{center}{\pgfpointorigin}
  \anchor{text}{\pgfpointorigin}
  \anchor{north}{\pgfpoint{0cm}{0.3cm}}
  \anchor{south}{\pgfpoint{0cm}{-0.3cm}}
  \anchor{east}{\pgfpoint{0.3cm}{0cm}}
  \anchor{west}{\pgfpoint{-0.3cm}{0cm}}
  \anchor{pin 1}{\pgfpoint{-0.7cm}{0cm}}
  \anchor{pin 2}{\pgfpoint{0cm}{-0.7cm}}
  \anchor{bpin 1}{\pgfpoint{-0.3cm}{0cm}}
  \anchor{bpin 2}{\pgfpoint{0cm}{-0.3cm}}
  \backgroundpath{
    \pgfpathmoveto{\pgfpoint{-0.274cm}{-0.122cm}}
    \pgfpatharc{204}{516}{0.3cm}
    \pgfpathmoveto{\pgfpoint{-0.7cm}{0cm}}\pgfpathlineto{\pgfpointorigin}
    \pgfpathmoveto{\pgfpoint{0cm}{-0.7cm}}\pgfpathlineto{\pgfpoint{0cm}{-0.3cm}}
  }
  \foregroundpath{
    \pgfpathcircle{\pgfpointorigin}{0.07cm}
    \pgfusepath{fill}
  }
}
\makeatother
\begin{document}
\begin{circuitikz}[american, line width=0.8pt]
\ctikzset{bipoles/length=1.2cm}
\ctikzset{grounds/scale=1.36}
\coordinate (b2) at (2,-2);
\coordinate (b5) at (8,-2);
\coordinate (b9) at (16,-2);
\coordinate (d5) at (8,-6);
\coordinate (f5) at (8,-10);
\coordinate (d9) at (16,-6);
\coordinate (f9) at (16,-10);
\coordinate (b12) at (22,-2);
\coordinate (c2) at (2,-4);
\coordinate (c12) at (22,-4);
\coordinate (g5) at (8,-12);
\coordinate (g9) at (16,-12);
\node[smacoax, draw, xscale=-1] (part-J1) at (b2) {}; % line 3
\node[anchor=south] at (part-J1.north) {$J_{1}$}; % line 3
\draw (b5) to[R, l_=$R_{1}$, a^=$18\,\Omega$] (b9); % line 4
\draw (d5) to[R, l_=$R_{2}$, a^=$300\,\Omega$] (f5); % line 5
\draw (d9) to[R, l_=$R_{3}$, a^=$300\,\Omega$] (f9); % line 6
\node[smacoax, draw] (part-J2) at (b12) {}; % line 7
\node[anchor=south] at (part-J2.north) {$J_{2}$}; % line 7
\node[ground] at (c2) {}; % line 8
\node[ground] at (c12) {}; % line 9
\node[ground] at (g5) {}; % line 10
\node[ground] at (g9) {}; % line 11
\draw (part-J1.pin 1) -- (b5); % line 13
\draw (b5) -- (d5); % line 14
\draw (b9) -- (part-J2.pin 1); % line 15
\draw (b9) -- (d9); % line 16
\draw (f5) -- (g5); % line 17
\draw (f9) -- (g9); % line 18
\draw (part-J1.pin 2) -- (c2); % line 19
\draw (part-J2.pin 2) -- (c12); % line 20
\node[circ] at (b5) {};
\node[circ] at (b9) {};
\node[anchor=south west, inner sep=2pt, circuitnotemark, font=\LARGE\bfseries] (circuittitle) at (current bounding box.north west) {X};
\path (circuittitle.south west) rectangle ++(12.619,0.853);
\end{circuitikz}
\end{document}
